3.6. Let $V$ be a locally free vector bundle of rank $n+2$ on $S$, and
let $X$ be the smooth quadric over $S$ defined by an ordinary quadratic
form $Q$ on $V$.  Let $H$ moreover be a hyperplane in
$\mathbb P(V^*)$ meeting $X$ transversally, and put $X^o=X-H$.  The
schemes $X$ and $Y=X\cap H$ are smooth quadrics over $S$, of dimensions
$n$ and $n-1$.  We may therefore use 3.3 to compute the cohomology of
$X^o=X-Y$.

%%%%%%%%%% source scan idx 86 | folio 79 %%%%%%%%%%

\begin{equation}\tag{3.6.1}
\begin{tikzcd}[column sep=huge,row sep=huge]
X^o \arrow[r,hook] \arrow[rd,"f"']
  & X \arrow[d,"p"]
  & Y \arrow[l,hook] \arrow[ld,"q"]\\
{} & S & \mathord{.}
\end{tikzcd}
\end{equation}

If $n=0$, then $Y=\varnothing$ and $X=X^o$.  Suppose that $n>0$.  We
have two long exact sequences dual to one another:
\begin{equation}\tag{3.6.2}
\cdots\xrightarrow{\partial}R^if_!\mathbb Z_\ell
\longrightarrow R^ip_*\mathbb Z_\ell
\xrightarrow{r_i}R^iq_*\mathbb Z_\ell\longrightarrow\cdots
\end{equation}
and
\begin{equation}\tag{3.6.3}
\cdots\xrightarrow{\partial}R^{i-2}q_*\mathbb Z_\ell(-1)
\longrightarrow R^ip_*\mathbb Z_\ell
\longrightarrow R^if_*\mathbb Z_\ell\longrightarrow\cdots.
\end{equation}

For $n$ even, the restriction homomorphism $r_i$ is an isomorphism for
$i\neq n,2n$; for $n$ odd, it is an isomorphism for $i\neq n-1$.  If
$n=2m$ and $\alpha$ is the cohomology class of a generator, then
\begin{equation*}
r_n(\alpha)=\tfrac12\eta^m,
\end{equation*}
so $r_n$ is surjective and its kernel is the primitive part of
$R^np_*\mathbb Z_\ell$.  If $n=2m+1$, then
$r_{2m}(\eta^m)=\eta^m$, so $r_{2m}$ is injective and its cokernel is the
primitive quotient of $R^{2m}q_*\mathbb Z_\ell$.  It follows that
$R^if_!\mathbb Z_\ell=0$ for $i\neq n,2n$, and that
$R^if_!\mathbb Z_\ell$ is a twisted constant free
$\mathbb Z_\ell$-sheaf of rank one for $i=n,2n$.  By duality, or from
(3.6.3), one likewise shows that $R^if_*\mathbb Z_\ell=0$ for
$i\neq0,n$, that $f_*\mathbb Z_\ell=\mathbb Z_\ell$, and that
$R^nf_*\mathbb Z_\ell$ has rank one.

\footnote{This paragraph prints $r=2m+1$ instead of $n=2m+1$ and cites
3.5.3 instead of the displayed sequence (3.6.3).}

Suppose that $n=2m$ and retain the notation of 3.3.

%%%%%%%%%% source scan idx 87 | folio 80 %%%%%%%%%%

\begin{center}
\begin{tikzcd}[column sep=large,row sep=large]
0 \arrow[r]
 & R^nf_!\mathbb Z_\ell(m) \arrow[r] \arrow[d,"\varphi"]
 & R^np_*\mathbb Z_\ell(m) \arrow[r] \arrow[d,equal]
 & R^nq_*\mathbb Z_\ell(m) \arrow[r]
 & 0\\
0
 & R^nf_*\mathbb Z_\ell(m) \arrow[l]
 & R^np_*\mathbb Z_\ell(m) \arrow[l]
 & R^{n-2}q_*\mathbb Z_\ell(m-1) \arrow[l]
 & 0\mathord{.} \arrow[l]
\end{tikzcd}
\end{center}

The sheaf $R^nf_!\mathbb Z_\ell(m)$ then has two natural opposite
generators $\pm\delta$, whose images in
$R^np_*\mathbb Z_\ell(m)$ are
$\pm(c\ell(\alpha)-c\ell(\beta))$.  By 3.3(iii)(b),
\begin{equation*}
\operatorname{Tr}(\delta^2)
=(c\ell(\alpha)-c\ell(\beta))^2=(-1)^m\cdot2,
\end{equation*}
and $\pm\varphi(\delta)$ is twice the natural dual generators
$\pm\delta'$ of $R^nf_*\mathbb Z_\ell(m)$, namely the opposite images
of $c\ell(\alpha)$ and $c\ell(\beta)$.

Suppose now that $n=2m+1$ and apply the notation of 3.3 to $Y$.

\begin{center}
\begin{tikzcd}[column sep=large,row sep=large]
0 \arrow[r]
 & R^{2m}p_*\mathbb Z_\ell(m) \arrow[r]
 & R^{2m}q_*\mathbb Z_\ell(m) \arrow[r,"\partial"]
 & R^nf_!\mathbb Z_\ell(m) \arrow[r] \arrow[d,"\varphi"]
 & 0\\
{}
 & R^{2m+2}p_*\mathbb Z_\ell(m+1)
 & R^{2m}q_*\mathbb Z_\ell(m) \arrow[l]
 & R^nf_*\mathbb Z_\ell(m+1) \arrow[l]
 & 0\mathord{.} \arrow[l]
\end{tikzcd}
\end{center}

\footnote{The leftmost term in the printed lower row has twist $(m)$;
duality and the other three terms require the displayed twist $(m+1)$.}

The sheaf $R^nf_!\mathbb Z_\ell(m)$ then has two natural opposite
generators $\pm\delta$, namely
$\partial c\ell(\alpha)$ and $\partial c\ell(\beta)$.  We have
\begin{equation*}
\delta\,\varphi(\delta)=\delta^2
=(\partial c\ell(\alpha))^2
=\partial\bigl(c\ell(\alpha)\,\partial c\ell(\alpha)\bigr)=0,
\end{equation*}
and therefore $\varphi(\delta)=0$.  On the other hand,
$R^nf_*\mathbb Z_\ell(m+1)$ has two natural opposite generators
$\pm\delta'$, whose images in $R^{2m}q_*\mathbb Z_\ell(m)$ are
$\pm(c\ell(\alpha)-c\ell(\beta))$, and
$\operatorname{Tr}(\delta\delta')=\pm1$.

%%%%%%%%%% source scan idx 88 | folio 81 %%%%%%%%%%

These results are assembled in the following table.

\underline{Table} 3.7. \underline{Cohomology of affine quadrics of
dimension $n$, $X^o=X-Y$.}

\medskip
\noindent\textbf{A.} The cohomology is torsion-free.  The nonzero Betti
numbers are
\begin{align*}
\text{ordinary cohomology:}\quad &b_0=b_n=1,\\
\text{cohomology with compact supports:}\quad &b_{2n}=b_n=1,
\end{align*}
except that $b_0=2$ when $n=0$.

\medskip
\noindent\textbf{B.} If $n=2m>0$, then
\begin{align*}
R^nf_!\mathbb Z_\ell(m)
  &=\text{the primitive part of }R^{2m}p_*\mathbb Z_\ell(m),\\
R^nf_*\mathbb Z_\ell(m)
  &=\text{the primitive quotient of }R^{2m}p_*\mathbb Z_\ell(m).
\end{align*}
If $n=2m+1$, then
\begin{align*}
R^nf_!\mathbb Z_\ell(m)
  &=\text{the primitive quotient of }R^{2m}q_*\mathbb Z_\ell(m),\\
R^nf_*\mathbb Z_\ell(m+1)
  &=\text{the primitive part of }R^{2m}q_*\mathbb Z_\ell(m).
\end{align*}
Locally, these groups have natural generators defined up to sign;
$\delta$ denotes those for cohomology with compact supports, and
$\delta'$ those for ordinary cohomology.

\medskip
\noindent\textbf{C.} The map
\begin{equation*}
\varphi:R^nf_!\mathbb Z_\ell\longrightarrow R^nf_*\mathbb Z_\ell
\end{equation*}
is zero when $n$ is odd, and sends $\pm\delta$ to $\pm2\delta'$ when
$n>0$ is even.  Moreover,
\begin{equation*}
\operatorname{Tr}(\delta\delta')=\pm1,
\qquad
\operatorname{Tr}(\delta^2)=
\begin{cases}
0, & n=2m+1,\\
(-1)^m\cdot2, & n=2m.
\end{cases}
\end{equation*}

\underline{Verification} 3.8. The complex affine quadric $S$ in
$\mathbb C^n$ with equation $\sum_i z_i^2=1$ is diffeomorphic to the
tangent bundle of the real sphere $S\cap\mathbb R^n$.  This agrees with
A and C.

\medskip
\underline{BIBLIOGRAPHY}

\begin{thebibliography}{9}
\bibitem{MumfordGIT}
D. Mumford, \emph{Geometric Invariant Theory}, Ergebnisse \textbf{34},
Springer-Verlag, 1965.
\end{thebibliography}
